\documentclass[a4paper,12pt]{article}
\usepackage[utf8]{inputenc}
\usepackage[T1]{fontenc}
\usepackage[french]{babel}
\usepackage[right=2.5cm, left=2.5cm]{geometry}
\usepackage[ddmmyyyy]{datetime}
\usepackage[table]{xcolor}
\usepackage{lmodern,mathptmx,changepage,titlesec,hyperref,listings,lstautogobble,graphicx,array,longtable,multirow,lipsum,tikz,shorttoc,enumitem}
\usetikzlibrary{arrows,automata}
\usetikzlibrary{positioning}

\renewcommand{\rmdefault}{\sfdefault} %Utilisation de la police sans-serif ("Computer Modern Sans") pour la police roman
\renewcommand{\ttdefault}{pcr} 	%Utilisation d'une police "CourrierNew" pour la police monospaced (pour faire un listing manuel)
\linespread{1.15}				%Interligne

%Utilisation de liens colorés en bleu et soulignés
\hypersetup{colorlinks=true, urlcolor=blue, urlbordercolor=blue, linkcolor=black, linkbordercolor=white}
\makeatletter \Hy@AtBeginDocument{\def\@pdfborder{0 0 1} \def\@pdfborderstyle{/S/U/W 1}}\makeatother

\titlespacing*{\section} {0cm}{7ex plus 1ex minus .2ex}{1.5ex plus .2ex}
\titlespacing*{\subsection} {0cm}{4.5ex plus 1ex minus .2ex}{1.5ex plus .2ex}
\titleformat*{\section}{\huge\bfseries}
\titleformat*{\subsection}{\Large\bfseries}
\titleformat*{\subsubsection}{\normalsize\bfseries}

\definecolor{darkgreen}{rgb}{0,0.8,0}
\definecolor{mygray}{rgb}{0.93,0.93,0.93}
\definecolor{mymauve}{rgb}{0.58,0,0.82}
\lstset{	
	basicstyle=\small\ttfamily,
	backgroundcolor=\color{mygray},
	breaklines=true,
	breakatwhitespace=true,
	postbreak=\raisebox{0ex}[0ex][0ex]{\ensuremath{\color{red}\hookrightarrow\space}},
	tabsize=3,
	frame=none,
	rulecolor=\color{black},
	keywordstyle=\color{blue}\bfseries,
	stringstyle=\color{orange},
	showstringspaces=false,
	commentstyle=\footnotesize\color{darkgreen},
	keepspaces=true,
	extendedchars=true,
	numbers=left,
	numberstyle=\tiny\color{lightgray},
	stepnumber=1,
	escapeinside={(@}{@)},
	autogobble=true,
	literate=
		{á}{{\'a}}1 {é}{{\'e}}1 {í}{{}}1 {ó}{{\'o}}1 {ú}{{\'u}}1
		{Á}{{\'A}}1 {É}{{\'E}}1 {Í}{{\'I}}1 {Ó}{{\'O}}1 {Ú}{{\'U}}1
		{à}{{\`a}}1 {è}{{\`e}}1 {ì}{{\`i}}1 {ò}{{\`o}}1 {ù}{{\`u}}1
		{À}{{\`A}}1 {È}{{\'E}}1 {Ì}{{\`I}}1 {Ò}{{\`O}}1 {Ù}{{\`U}}1
		{ä}{{\"a}}1 {ë}{{\"e}}1 {ï}{{\"i}}1 {ö}{{\"o}}1 {ü}{{\"u}}1
		{Ä}{{\"A}}1 {Ë}{{\"E}}1 {Ï}{{\"I}}1 {Ö}{{\"O}}1 {Ü}{{\"U}}1
		{â}{{\^a}}1 {ê}{{\^e}}1 {î}{{\^i}}1 {ô}{{\^o}}1 {û}{{\^u}}1
		{Â}{{\^A}}1 {Ê}{{\^E}}1 {Î}{{\^I}}1 {Ô}{{\^O}}1 {Û}{{\^U}}1
		{œ}{{\oe}}1 {Œ}{{\OE}}1 {æ}{{\ae}}1 {Æ}{{\AE}}1 {ß}{{\ss}}1
		{ç}{{\c c}}1 {Ç}{{\c C}}1 {ø}{{\o}}1 {å}{{\r a}}1 {Å}{{\r A}}1
		{€}{{e}}1 {£}{{\pounds}}1 {«}{{\guillemotleft}}1
		{»}{{\guillemotright}}1 {ñ}{{\~n}}1 {Ñ}{{\~N}}1 {¿}{{?`}}1
}

%Redéfinition de la taille de \Huge pour le titre du document
\makeatletter\renewcommand\Huge{\@setfontsize\Huge{37pt}{40}}\makeatother
\date{\today}

\title{\vspace{\fill}\textbf{\Huge Compte Rendu}}
\author{
	Sonny Klotz - Jean-Didier Pailleux - Malek Zemni
	\vspace{2em}\\
	\textit{Interface de chargement, de contrôle}\\\textit{et d’analyse statistique des données}\\\textit{pour la constitution d’un graphe de flux}
	\vspace{2em}
}

\begin{document}
\pagenumbering{gobble}\clearpage
\maketitle\vspace{9em}
\begin{center}\includegraphics[scale=0.7]{../Cahier/logo.png}\end{center}
\begin{flushright}Module \textit{Projet}\end{flushright}
\newpage
\tableofcontents
\newpage\clearpage\pagenumbering{arabic}

	\section*{Introduction}
		Motivations du projet : un mix de
			- Description Sujet (notre premier oral)
			- Intro et partie 1 cahier charges
		
		Plan :
			- archi appli (cahier charges) + données transitent + fonctionnalités
			- langages de programmation : découlant des contraintes de l'appli
			- bilan technique
			- organisation interne
			
		Documents fournis en plus du code :
			- Manuel
			- Documentation
			
		-> 1 page (Sonny)
	
	\section{Architecture de l'application}
		\subsection{Organigramme et données échangées}
			Phrase de présentation -> diagramme -> légende (non modifiée, on la met en accord avec le dernier produit quand on aura finit)
		\subsection{Format du fichier CSV}
			une demi-page
		\subsection{Fonctionnalités des modules}
			Les fonctionnalités de notre produit final (pas ceux du cahier des charges)
		
		-> entre 3 pages (pas plus) pour la partie (Sonny)
		
	\section{Contraintes et Langages de programmation}
		contraintes :
			-voir cahier des charges : applet, add, ...
		choix langages et frameworks:
			-python : justification + flask (voir cahier + ton expérience) -> intègre en plus un module de tests
			-html / css : easy on fait du web
			-javascript - jquery et AJAX: UX design tu peux parler vite fait, js donne intéractivité coté client, ajax pour communiquer infos côté serveur entre autres et jquery qui permet de faire tout ca
			-Sphinx, génération documentation automatique (la doc flask est générée grace a sphinx par ex), si tu sais pas on le fera plus tard, mais faut pas l'oublier
			
		-> 2 pages (J-D)
		
	\section{Partie technique du projet}
		Quoi marche (comment ca marche - à l'aide de la partie avant)
		QUoi marche pas (pb et et resolution si trouvee)
		
		1 page, plus tard
		
	\section{Points délicats}
		
		Justification des points qui divergent par rapport aux specs et organigramme de départ
		
		1 page, plus tard à moins que t'ai déjà des idées
				
	\section{Organisation interne du groupe}
		
		Insérer la répartition des tâches (cahier des charges)
		
		-> (J-D)
		
		1 page
		
		On expliquera plus tard comment ca s'est déroulé concrètement par rapport à ce quiétait prévu.
		
	\section{Coût}
		
		Insérer le tableau des couts, ajouter une colonne (vide) "coût" le final
		
		-> J-D
		1 demi - page
		
	\section*{Conclusion}
		
		1 demi-page, plus tard
		
\end{document}
